% source: https://github.com/pfradin/luadraw/discussions/320#discussioncomment-18219103 (pfradin, block 1)
\documentclass{standalone}
\usepackage[3d]{luadraw}
\usepackage[svgnames]{xcolor}

% Color palette definitions (Ocean Blue + Light Green Cartographic Theme)
\definecolor{earthocean}{RGB}{148,196,237}   % Ocean Blue
\definecolor{earthland}{RGB}{196,226,182}    % Pastel Cartographic Light Green
\definecolor{earthcoast}{RGB}{110,165,210}   % Crisp Shoreline Blue
\definecolor{earthedge}{RGB}{110,165,210}    % Spherical Horizon Outline
\definecolor{earthaxis}{RGB}{160,45,45}      % Earth Axis Burgundy

\begin{document}
\begin{luadraw}{name=globe_110m_simple}
    local ld = luadraw
    require "luadraw_spherical"

    -- Initialize 3D viewport and spherical graph
    local g = ld.graph3d:new{
        window3d = {-4, 4, -4, 4, -4, 4},
        window = {-4, 4, -4, 4},
        viewdir = {30, 65},
        size = {12, 12}
    }

    -- Helper to load Lua dataset with dual path resolution
    local function load_dataset(filename)
        local path = "generated/" .. filename
        local f = io.open(path, "r")
        if f then
            f:close()
            return dofile(path)
        end
        return dofile("demo/generated/" .. filename)
    end

    -- 1. Load 3 core 110m datasets
    local ocean = load_dataset("ocean_110m.lua")
    local land = load_dataset("land_110m.lua")
    local coastline = load_dataset("coastline_110m.lua")

    -- Axial tilt (Obliquity of the ecliptic: 23.44 deg)
    local axial_tilt = -23.439281

    -- Helper to project (lon, lat) points onto the 3D sphere surface with axial tilt
    local function convert(points)
        local result = {}
        for _, point in ipairs(points) do
            table.insert(result, ld.sM(point.lon, 90 - point.lat)) --<- change here
        end
        return ld.rotate3d(result, axial_tilt, {ld.pt3d.Origin, ld.pt3d.vecJ})
    end

    -- Base Sphere definition
    g:Define_sphere{
        color = "earthocean",
        edgecolor = "earthedge",
        edgewidth = 2.5,
        back = false, inside = false --<- for version 3.5
    }

    -- Draw tilted Earth axis (passing through center and protruding slightly beyond poles)
    local axis_direction = ld.rotate3d(ld.pt3d.vecK, axial_tilt, {ld.pt3d.Origin, ld.pt3d.vecJ})
    local axis_path = {-4.1 * axis_direction, 4.1 * axis_direction}
    g:Dseg3d(axis_path, "earthaxis,line width=1.2") --<- drawn behind the sphere
    g:DSseg(axis_path, {color="earthaxis", width=12}) --<- with this, only the front part is visible

    -- Layer 1: Ocean Polygons
    for _, feature in ipairs(ocean) do
        for _, ring in ipairs(feature.rings) do
            g:DSregion(convert(ring.points), {
                style = "noline",
                fill = "earthocean",
                fillopacity = 1
            })
        end
    end

    -- Layer 2a: Land Outer Rings (Continents and Islands)
    for _, feature in ipairs(land) do
        for _, ring in ipairs(feature.rings) do
            if ring.role == "outer" then
                g:DSregion(convert(ring.points), {
                    style = "noline",
                    fill = "earthland",
                    fillopacity = 1
                })
            end
        end
    end

    -- Layer 2b: Land Inner Rings (Inland holes / Caspian Sea backfill)
    for _, feature in ipairs(land) do
        for _, ring in ipairs(feature.rings) do
            if ring.role == "inner" then
                g:DSregion(convert(ring.points), {
                    style = "noline",
                    fill = "earthocean",
                    fillopacity = 1
                })
            end
        end
    end

    -- Layer 3: Coastline Curves
    for _, feature in ipairs(coastline) do
        g:DScurve(convert(feature.points), {
            color = "earthcoast",
            width = 0.55
        })
    end

    -- Compute 3D spherical visibility and occlusion, then emit TikZ code
    g:Dspherical()
    g:Show()
\end{luadraw}
\end{document}
